%!TEX root = ./main.tex

\section[Allocation]{Choosing a Point on the Curve}

% SECTION TIMING: 20:00--32:00 (5 frames; the peak explainer is a ~1 min beat).
% Questions 2 and 3 of the roadmap frame: the opening frame answers Q2 in one slide
% now that the OFDM primer has moved to Class 8, and the rest answers Q3. The simulation-parameters frame and the multiuser
% ODDM frame were removed on 2026-09-13; see parked/05_simulation_multiuser_frames.tex.
% SOURCE: "Cross-Domain Dual-Functional OFDM Waveform Design for Accurate
% Sensing/Positioning," IEEE JSAC, 2024. The two allocation strategies and their
% "almost no loss" results are quoted from that work.

% Moved here 2026-09-20 from 03_waveforms.tex, which was parked whole. This is now
% the lecture's entire answer to Question 2 ("which signal do we actually use?") and
% the setup for Question 3: it names OFDM, says why it costs nothing to adopt, and
% turns the grid into the thing we are about to hand out. The ground-up primer that
% used to precede it is taught in Class 8 (Wi-Fi History) instead.
\begin{frame}{OFDM already carries a radar inside it}
  \vspace{0.3em}
  \begin{columns}[T,onlytextwidth]
    \column{0.44\textwidth}
      \centering
      \begin{tikzpicture}[scale=0.80,transform shape]
        \foreach \x in {0,...,5}
          \foreach \y in {0,...,3}
            \draw[draw=softgray!50,fill=blue!6] (\x*0.62+0.12,\y*0.52+0.10)
              rectangle ++(0.58,0.48);
        \draw[-{Latex[length=2mm]},softgray] (0,0) -- (4.0,0)
          node[anchor=north,font=\scriptsize] {time (symbols)};
        \draw[-{Latex[length=2mm]},softgray] (0,0) -- (0,2.4);
        \node[font=\scriptsize,rotate=90,anchor=south] at (-0.16,1.2) {frequency};
      \end{tikzpicture}
    \column{0.52\textwidth}
      \small
      \begin{itemize}\itemsep5pt
        \item Mature: Wi-Fi, 4G, 5G, 6G.
        \item The grid is already knobs.
        \item A cell serves bits or echoes.
        \item No new modulator needed.
      \end{itemize}
  \end{columns}

  \vspace{0.4em}
  \qa{If OFDM is already good enough, what is left to design?}
     {Which cells go to bits and which go to echoes. Every cell handed to sensing
      is a cell taken from data.}

  \glossline{OFDM = Orthogonal Frequency Division Multiplexing}
\end{frame}

\begin{frame}{The grid is the budget}
  \vspace{0.3em}
  \begin{columns}[T,onlytextwidth]
    \column{0.50\textwidth}
      \centering
      \begin{tikzpicture}[scale=0.78,transform shape]
        \foreach \x/\y in {0/0,1/0,3/0,4/0,1/1,2/1,5/1,0/2,3/2,4/2,2/3,5/3,0/3,3/3}
          \fill[blue!22,draw=softgray!50] (\x*0.62+0.14,\y*0.52+0.12) rectangle ++(0.58,0.48);
        \foreach \x/\y in {2/0,5/0,0/1,3/1,4/1,1/2,2/2,5/2,1/3,4/3}
          \fill[orange!35,draw=softgray!50] (\x*0.62+0.14,\y*0.52+0.12) rectangle ++(0.58,0.48);
        \draw[-{Latex[length=2mm]},softgray] (0,0) -- (4.05,0)
          node[anchor=north,font=\scriptsize] {time};
        \draw[-{Latex[length=2mm]},softgray] (0,0) -- (0,2.42);
        \node[font=\scriptsize,rotate=90,anchor=south] at (-0.16,1.2) {frequency};
      \end{tikzpicture}

      \vspace{0.4em}
      {\scriptsize\textcolor{ranblue}{$\blacksquare$ communication RE} \quad
       \textcolor{coreorange}{$\blacksquare$ sensing RE}\par}
    \column{0.46\textwidth}
      \small
      \begin{itemize}\itemsep5pt
        \item One RE is one cell.
        \item Each carries symbol and power.
        \item Assign cell, then power.
        \item Two functions, one grid.
      \end{itemize}
  \end{columns}

  \glossline{RE = Resource Element, one subcarrier during one OFDM symbol\\
    PSLR = Peak-to-SideLobe Ratio; RoI = Region of Interest: the distances and speeds we actually care about}
  \source{https://ieeexplore.ieee.org/document/10683983}{Cross-Domain Dual-Functional OFDM Waveform Design, IEEE JSAC 2024}
\end{frame}

% Inserted 2026-09-09. PSLR grades the next three frames, but "peak", "sidelobe"
% and "clutter" were never defined -- students do not know what the peak of a
% signal is. Define all three with one picture before anything gets graded.
\begin{frame}{What a peak is, before we grade one}
  \vspace{0.2em}
  \centering
  \begin{tikzpicture}[scale=0.94,transform shape]
    % axes: matched-filter output against delay
    \draw[-{Latex[length=2.2mm]},softgray] (0,0) -- (8.6,0)
      node[anchor=north,font=\scriptsize,xshift=-9mm] {delay (= distance)};
    \draw[-{Latex[length=2.2mm]},softgray] (0,0) -- (0,3.0);
    \node[font=\scriptsize,rotate=90,anchor=south] at (-0.2,1.4) {receiver output};
    % clutter floor
    \draw[thick,draw=softgray!80]
      (0.15,0.16) .. controls (0.7,0.34) and (1.0,0.08) .. (1.5,0.22)
      .. controls (1.9,0.34) .. (2.3,0.14)
      -- (2.3,0.14)
      (5.6,0.20) .. controls (6.1,0.10) and (6.6,0.36) .. (7.1,0.16)
      .. controls (7.6,0.30) .. (8.2,0.18);
    % left sidelobe, main peak, right sidelobe
    \draw[very thick,draw=ranblue]
      (2.3,0.14) .. controls (2.55,0.14) .. (2.75,0.62)
      .. controls (2.95,0.14) and (3.15,0.14) .. (3.35,0.14)
      .. controls (3.7,0.2) .. (3.95,2.45)
      .. controls (4.2,0.2) and (4.4,0.14) .. (4.65,0.14)
      .. controls (4.85,0.14) .. (5.05,0.62)
      .. controls (5.25,0.14) .. (5.6,0.20);
    % labels
    \node[font=\scriptsize\bfseries,text=ranblue,anchor=south] at (3.95,2.5) {the peak: a target};
    \node[font=\scriptsize,text=ranblue,anchor=south] at (2.45,0.66) {sidelobes};
    \node[font=\scriptsize,text=softgray,anchor=south] at (1.1,0.4) {clutter};
    % PSLR arrow
    \draw[{Latex[length=1.8mm]}-{Latex[length=1.8mm]},coreorange,thick]
      (6.4,0.62) -- node[font=\scriptsize\bfseries,text=coreorange,right=2pt]{PSLR}
      (6.4,2.45);
    \draw[dotted,coreorange] (3.95,2.45) -- (6.4,2.45);
    \draw[dotted,coreorange] (5.05,0.62) -- (6.4,0.62);
  \end{tikzpicture}

  \vspace{0.45em}
  {\small The receiver slides a copy of its own waveform across the echo.
   Where a target sits, everything lines up and adds --- that spike is
   \textbf{the peak}. \textbf{Sidelobes} are ghosts of our own waveform;
   \textbf{clutter} is every echo we never asked for.\par}

  \vspace{0.35em}
  \takeaway{PSLR = how far the true peak stands above the tallest ghost.\\That number is the sensing grade on the next slide.}
\end{frame}

\begin{frame}{Two scores, graded in two rooms}
  \vspace{0.6em}
  \begin{columns}[T,onlytextwidth]
    \column{0.47\textwidth}
      \large\textbf{\textcolor{ranblue}{Communication score}}

      \vspace{0.5em}
      \small
      Achievable data rate, summed over every RE assigned to communication.

      \vspace{0.5em}
      Measured in \textbf{time--frequency}: put bits where the channel is good.
    \column{0.47\textwidth}
      \large\textbf{\textcolor{coreorange}{Sensing score}}

      \vspace{0.5em}
      \small
      PSLR inside the region of interest --- how far the true peak stands above the clutter.

      \vspace{0.5em}
      Measured in \textbf{distance and speed}: make the peak stand tall where the
      targets actually are.
  \end{columns}

  \vspace{0.8em}
  \takeaway{The same grid is graded twice,\\by two examiners in two different domains.}
\end{frame}

\begin{frame}{Whoever picks first, wins}
  \vspace{0.4em}
  \begin{columns}[T,onlytextwidth]
    \column{0.47\textwidth}
      {\small\textbf{\textcolor{ranblue}{Communication-centric}}\par}
      \vspace{0.4em}
      \small
      \begin{enumerate}\itemsep4pt
        \item Give the good REs to bits.
        \item Then tune leftover powers.
        \item Sensing takes the remainder.
      \end{enumerate}

      \vspace{0.4em}
      {\scriptsize\textcolor{softgray}{Rate stays optimal; PSLR is merely good inside the RoI.}\par}
    \column{0.47\textwidth}
      {\small\textbf{\textcolor{coreorange}{Sensing-centric}}\par}
      \vspace{0.4em}
      \small
      \begin{enumerate}\itemsep4pt
        \item Lock the peak shape inside the RoI.
        \item Then tune the rest of the grid.
        \item Bits take the remainder.
      \end{enumerate}

      \vspace{0.4em}
      {\scriptsize\textcolor{softgray}{PSLR is locally perfect; rate approaches, not reaches, optimal.}\par}
  \end{columns}

  \vspace{0.7em}
  \qa{Both designs claim ``almost no loss''. How can that be?}
     {Because the constraint is only enforced inside the \textbf{region of
      interest}. Narrow the room you must be perfect in, and the rest of the grid
      is free for the other function.}
\end{frame}
